\section{Data Preprocessing and Treatment}

Data preprocessing was conducted to ensure data quality and consistency while maximizing retention of all valid observations. The preprocessing pipeline consisted of five primary stages: data integration, standardization, missing data handling, quality assurance, and final dataset construction.

\subsection{Data Integration and Standardization}

Individual participant data files (\textit{N} = 170) generated by PsychoPy were systematically merged to create a comprehensive trial-level dataset. Of these, 170 files contained valid experimental data and were retained for analysis. Demographic information (age, gender, handedness, and vision status) was integrated from a separate data file through participant ID matching. Key response variables including accuracy (\texttt{resp.corr}), response time (\texttt{resp.rt}), and confidence ratings (\texttt{conf\_radio.response}) were converted to numeric format to ensure computational consistency. Participant identification codes were standardized across datasets through case normalization and whitespace trimming to facilitate accurate data linking across multiple files.

\subsection{Missing Data Treatment}

A critical decision was made to adopt a \textit{no-data-loss} approach for participant retention. All 146 participants who completed the recognition memory task were retained in the final sample, regardless of demographic data completeness. Among these participants, 38 individuals (26.0\%) had missing demographic information in one or more variables. Rather than employing listwise deletion, which would have reduced statistical power and potentially introduced selection bias, a condition-wise imputation strategy was implemented. For the continuous variable age, missing values were imputed using the median value within each experimental condition (Natural Boundary vs. Abrupt Boundary). For categorical variables (gender, handedness, vision), missing values were replaced with the mode within each respective condition. This condition-stratified approach preserved potential group differences while avoiding cross-condition information leakage.

To maintain analytical transparency and enable sensitivity analyses, two sets of indicator variables were created for each demographic measure. Binary flags (\texttt{*\_was\_missing}) identified originally missing values, while imputation indicators (\texttt{*\_imputed}) marked successfully imputed cases. This dual-flagging system allows researchers to assess whether patterns of missingness or the imputation procedure itself influenced analytical outcomes.

\subsection{Quality Assurance Procedures}

Systematic data quality checks were performed to identify potential data anomalies or recording errors. Recognition accuracy values were verified to fall within the valid probability range [0, 1], with all participant means conforming to this constraint. Response time distributions were examined for extreme outliers using the interquartile range (IQR) method with a conservative threshold of 3×IQR beyond the first and third quartiles. Identified outliers were flagged with a binary indicator (\texttt{rt\_outlier}) but retained in the dataset to preserve the full range of response patterns. Confidence ratings were confirmed to fall within the expected Likert scale range of 1 to 5. Trial counts per participant were validated against the expected value of 40 trials, with any deviations noted for further investigation.

\subsection{Final Dataset Structure}

The preprocessing pipeline yielded three complementary datasets optimized for different analytical approaches. The participant-level dataset contains 146 rows (79 Natural Boundary, 67 Abrupt Boundary) with aggregated performance metrics including mean accuracy, median response time, response time variability, and mean confidence rating for each individual. The trial-level dataset preserves all 6,800 individual trial responses, enabling fine-grained analysis of within-subject variability and trial-specific effects. A merged dataset combines trial-level responses with participant demographics, facilitating multilevel modeling approaches that account for both within-subject and between-subject sources of variation. This hierarchical data structure supports both traditional ANOVA-style analyses at the participant level and more sophisticated mixed-effects modeling at the trial level.

The preprocessing approach prioritized methodological transparency and reproducibility. All data transformations were documented with explicit flags, and no valid experimental data were excluded based on demographic missingness. This conservative approach maximizes statistical power while providing the analytical flexibility to assess the robustness of findings across different data treatment decisions.
